\section{Experiments}\label{sec:experiments}

\subsection{Experimental Setup}

\paragraph{Tasks.}
In \emph{simulation} (MuJoCo), a robot strikes a puck toward a target near a table edge or pushes a box close to an edge, with three surfaces (friction \ph{0.60}, \ph{0.35}, \ph{0.12}) and three loads (\ph{0.2}, \ph{0.6}, \ph{1.0}\,kg); the world model is a DINO-WM-style predictor on frozen DINOv2 features~\citep{zhou2025dinowm} trained on a single condition. On the \emph{benchmarks}, we use PushT and PointMaze with the released AdaJEPA checkpoints~\citep{wang2026adajepa} and pair the existing agent and block recolourings with changes in mass and damping, so that colour becomes an imperfect cue. On \emph{hardware}, a Franka Panda strikes a puck on interchangeable mats and lifts containers of unknown mass using a frozen V-JEPA~2-AC predictor~\citep{assran2025vjepa2}. During a \ph{1{,}200}-episode deployment, lab members exchanged mats and containers on a schedule unknown to the robot.

\paragraph{Protocol and baselines.}
Each stream has $N=\ph{500}$ episodes over $K=9$ surface--load contexts with geometric dwell times. The cue shows the true context's appearance with probability $\rho$ and each other appearance with probability $(1-\rho)/(K-1)$, so $\rho=1/K$ is uninformative. We report 95\% bootstrap intervals over \ph{10} streams shared by all methods, and compare against the \emph{Frozen} world model; \emph{Reset-SGD}, which takes AdaJEPA-style gradient steps and resets every episode~\citep{wang2026adajepa}; \emph{Reset-Residual}, a Sandwich-Residuals-style adapter that also resets~\citep{soni2026sandwich}; \emph{Continual-SGD}, the same updates without resets; \emph{Mixture}, a persistent mixture of dynamics models selected from prediction error alone~\citep{nagabandi2019mole}; \emph{\method-NoCue}, our method without cues; and \emph{Oracle}, which uses \method's memories with the true context. All methods share backbone, planner, candidates and action budget; hyperparameters are tuned on separate development streams.

% PLACEHOLDER DATA: every coordinate in this file is illustrative and must be
% replaced by measured values before submission.
\begin{figure}[t]
\centering
\begin{tikzpicture}
\begin{groupplot}[
  group style={group size=4 by 1, horizontal sep=0.95cm},
  width=0.29\linewidth, height=3.3cm,
  tick label style={font=\tiny}, label style={font=\scriptsize}, title style={font=\scriptsize, yshift=-3pt},
  ylabel style={yshift=-3pt},
  legend style={font=\tiny, draw=none, fill=none},
  every axis plot/.append style={thick, mark size=1.2pt},
]
% ---------------------------------------------------------------- (a) error vs cue reliability
\nextgroupplot[
  title={(a) first-contact error},
  xlabel={cue reliability $\rho$}, ylabel={error / frozen},
  xmin=0.08, xmax=1.0, ymin=0, ymax=1.12,
  xtick={0.11,0.5,0.8,0.95}, xticklabels={$\frac1K$,.5,.8,.95},
  legend to name=leg:results, legend columns=6,
]
\addplot[cFrozen, dashed] coordinates {(0.11,1.0) (0.95,1.0)}; \addlegendentry{Frozen}
\addplot[cReset, mark=x, only marks] coordinates {(0.11,1.0) (0.35,1.0) (0.65,1.0) (0.8,1.0) (0.95,1.0)}; \addlegendentry{Reset-SGD / Residual}
\addplot[cCont, mark=diamond*] coordinates {(0.11,0.86) (0.35,0.85) (0.65,0.87) (0.8,0.85) (0.95,0.86)}; \addlegendentry{Continual-SGD}
\addplot[cNoCue, mark=square*] coordinates {(0.11,0.63) (0.35,0.62) (0.65,0.63) (0.8,0.62) (0.95,0.62)}; \addlegendentry{Mixture / \method-NoCue}
\addplot[name path=up, draw=none, forget plot] coordinates {(0.11,0.68) (0.35,0.55) (0.65,0.39) (0.8,0.29) (0.95,0.18)};
\addplot[name path=lo, draw=none, forget plot] coordinates {(0.11,0.58) (0.35,0.46) (0.65,0.31) (0.8,0.21) (0.95,0.12)};
\addplot[cCira!25, forget plot] fill between[of=up and lo];
\addplot[cCira, mark=*] coordinates {(0.11,0.63) (0.35,0.50) (0.65,0.35) (0.8,0.25) (0.95,0.15)}; \addlegendentry{\method}
\addplot[cOracle, densely dotted] coordinates {(0.11,0.09) (0.95,0.09)}; \addlegendentry{Oracle}
\node[anchor=east, font=\tiny, text=red!70!black] at (rel axis cs:0.98,0.36) {placeholder};

% ---------------------------------------------------------------- (b) recall vs learning
\nextgroupplot[
  title={(b) recall vs.\ learning},
  ybar stacked, bar width=6pt,
  xlabel={occurrence}, ylabel={gain / initial loss},
  ymin=0, ymax=1.35, ytick={0,0.4,0.8}, xtick={1,2,3,4}, xticklabels={1st,2nd,3rd,4th+},
  enlarge x limits=0.18,
  legend style={at={(0.5,0.98)}, anchor=north, font=\tiny, legend columns=2},
]
\addplot[fill=cCira, draw=cCira] coordinates {(1,0.04) (2,0.52) (3,0.63) (4,0.70)}; \addlegendentry{recall $R$}
\addplot[fill=cCira!30, draw=cCira] coordinates {(1,0.78) (2,0.26) (3,0.17) (4,0.12)}; \addlegendentry{learning $P$}
\node[anchor=north, font=\tiny, text=red!70!black] at (rel axis cs:0.5,0.84) {placeholder};

% ---------------------------------------------------------------- (c) cumulative first-contact error
\nextgroupplot[
  title={(c) cumulative first-contact bias},
  xlabel={episode}, ylabel={$\sum_e b_e$ (norm.)},
  xmin=0, xmax=500, ymin=0, ymax=330,
  scaled x ticks=false, xtick={0,250,500},
  legend style={at={(0.02,0.98)}, anchor=north west, font=\tiny, fill=white, fill opacity=0.8, text opacity=1, row sep=-2pt},
]
\addplot[cReset, domain=0:500, samples=20] {0.62*x}; \addlegendentry{reset}
\addplot[cNoCue, domain=0:500, samples=40] {0.42*x - 10*ln(1+x/40)}; \addlegendentry{no cue}
\addplot[cCira, domain=0:500, samples=40] {35*ln(1+x/12) + 0.03*x}; \addlegendentry{\method}
\addplot[cCira, dashed, domain=0:500, samples=40] {52*ln(1+x/12) + 0.08*x}; \addlegendentry{bound}
\node[anchor=south east, font=\tiny, text=red!70!black] at (rel axis cs:0.98,0.03) {placeholder};

% ---------------------------------------------------------------- (d) unseen combinations
\nextgroupplot[
  title={(d) unseen combinations},
  ybar, bar width=4pt, ymin=0, ymax=1.45, ytick={0,0.5,1},
  ylabel={held-out error / frozen},
  symbolic x coords={conn,disc}, xtick=data, xticklabels={connected,disconn.},
  enlarge x limits=0.5,
  legend style={at={(0.5,0.99)}, anchor=north, font=\tiny, legend columns=2, /tikz/every even column/.append style={column sep=2pt}},
]
\addplot[fill=cFrozen!50, draw=cFrozen] coordinates {(conn,0.92) (disc,0.93)}; \addlegendentry{indep.}
\addplot[fill=cNoCue!60, draw=cNoCue] coordinates {(conn,0.41) (disc,0.88)}; \addlegendentry{$\kappa_\times{=}0$}
\addplot[fill=cCira, draw=cCira] coordinates {(conn,0.33) (disc,0.84)}; \addlegendentry{\method}
\addplot[fill=cOracle!60, draw=cOracle] coordinates {(conn,0.24) (disc,0.26)}; \addlegendentry{oracle}
\end{groupplot}
\end{tikzpicture}
\\[1pt]
\pgfplotslegendfromname{leg:results}
\caption{\textbf{Main results in simulation} (placeholder data; bands and error bars are 95\% intervals over \ph{10} streams). (a)~First-contact prediction error relative to the frozen model. Reset methods coincide with the frozen model at every cue reliability, whereas \method improves as cues become informative. (b)~Decomposition of the improvement into recall $R$ and learning $P$ (\cref{eq:decomp}) by how often a context has occurred. (c)~Cumulative first-contact bias grows linearly for reset methods and sublinearly for \method, below the bound of \cref{prop:first}. (d)~Error at the first contact with held-out surface--load combinations; the factored prior transfers only when the observed combinations are connected (\cref{prop:composition}, design in \cref{fig:composition}).}
\label{fig:results}
\end{figure}


\subsection{Results}

\paragraph{\method improves the first contact when contexts recur.}
\Cref{fig:results}a shows first-contact prediction error in simulated striking as a function of cue reliability. As predicted by \cref{obs:first}, Reset-SGD and Reset-Residual coincide with the frozen model. History alone (Mixture, \method-NoCue) reduces the error by \ph{XX\%}, and informative cues reduce it to \ph{XX\%} below the best reset method at $\rho=0.8$, close to the oracle at $\rho=0.95$. Uninformative cues do no harm. The cumulative first-contact bias in \cref{fig:results}c grows linearly for reset methods and sublinearly for \method, and it stays below the bound of \cref{prop:first} evaluated with the recorded beliefs. Without recurrence, \method provides no benefit (\ph{$\pm$X\%}), as the bound predicts. The gains carry over to task success on all platforms (\cref{tab:main}); on hardware, strike success rises from \ph{38\%} to \ph{74\%}.

% PLACEHOLDER DATA: every number marked \ph{} is unmeasured.
\begin{table}[t]
\centering
\scriptsize
\setlength{\tabcolsep}{3.2pt}
\resizebox{\linewidth}{!}{%
\begin{tabular}{@{}l cc cc cc c@{}}
\toprule
& \multicolumn{2}{c}{Sim.\ strike ($\rho{=}0.8$)} & \multicolumn{2}{c}{Benchmarks: first-contact err.\,$\downarrow$} & \multicolumn{2}{c}{Hardware strike} & \\
\cmidrule(lr){2-3}\cmidrule(lr){4-5}\cmidrule(lr){6-7}
Method & FC err.\,$\downarrow$ & success\,$\uparrow$ & PushT & PointMaze & stop err.\ [cm]\,$\downarrow$ & success\,$\uparrow$ & ms/step \\
\midrule
Frozen                                        & \ph{1.00} & \ph{41.2} & \ph{1.00} & \ph{1.00} & \ph{14.8} & \ph{38} & \ph{21} \\
Reset-SGD~\citep{wang2026adajepa}             & \ph{1.00} & \ph{41.6} & \ph{1.00} & \ph{1.00} & \ph{14.9} & \ph{39} & \ph{58} \\
Reset-Residual~\citep{soni2026sandwich}       & \ph{1.00} & \ph{41.9} & \ph{1.00} & \ph{1.00} & \ph{14.6} & \ph{40} & \ph{34} \\
Continual-SGD                                 & \ph{0.85} & \ph{47.3} & \ph{0.88} & \ph{0.83} & \ph{12.7} & \ph{45} & \ph{58} \\
Mixture~\citep{nagabandi2019mole}             & \ph{0.63} & \ph{58.0} & \ph{0.69} & \ph{0.61} & \ph{9.8}  & \ph{55} & \ph{29} \\
\method-NoCue                                 & \ph{0.62} & \ph{58.4} & \ph{0.67} & \ph{0.60} & \ph{9.5}  & \ph{56} & \ph{24} \\
\method                                       & \textbf{\ph{0.25}} & \textbf{\ph{76.5}} & \textbf{\ph{0.41}} & \textbf{\ph{0.32}} & \textbf{\ph{5.1}} & \textbf{\ph{74}} & \ph{24} \\
\midrule
\textcolor{black!60}{Oracle context}          & \textcolor{black!60}{\ph{0.09}} & \textcolor{black!60}{\ph{83.1}} & \textcolor{black!60}{\ph{0.22}} & \textcolor{black!60}{\ph{0.17}} & \textcolor{black!60}{\ph{3.6}} & \textcolor{black!60}{\ph{81}} & \textcolor{black!60}{\ph{24}} \\
\bottomrule
\end{tabular}}
\caption{\textbf{Main results} (placeholder values). First-contact (FC) error is normalized by the frozen model; success means stopping within \ph{5}\,cm of the target without leaving the table. Hardware results exclude first occurrences of each condition.}
\label{tab:main}
\end{table}


\paragraph{The gain on recurring contexts is recall.}\label{sec:decomp}
To attribute improvements, we record the adaptive state $(\pi,\Theta)$ before the episode's cue (index $0$) and after $s$ post-contact observations (index $1$), and evaluate all four combinations $L_{ij}=L(\pi^{(i)},\Theta^{(j)})$ on held-out transitions of the current context. The total gain $G=L_{00}-L_{11}$ then decomposes into the two-player Shapley values
\begin{equation}\label{eq:decomp}
  R=\tfrac12\big[(L_{00}-L_{10})+(L_{01}-L_{11})\big],\qquad P=\tfrac12\big[(L_{00}-L_{01})+(L_{10}-L_{11})\big],
\end{equation}
where $R$ corresponds to apparent learning, or recall, and $P$ to proper learning~\citep{heald2021coin}; reset methods have $R=0$ by construction. On the first occurrence of a context, nearly all of \method's gain is learning, whereas from the second occurrence onward recall accounts for \ph{XX\%} (\cref{fig:results}b). Freezing all memory parameters after calibration raises the error on returns by only \ph{XX\%}.

\paragraph{The factored prior transfers only along connected combinations.}
We train on six of the nine surface--load combinations and evaluate the first contact with each of the remaining three (design in \cref{fig:composition}, \cref{app:additional}). When the observed combinations connect all factor values, the factored prior reduces held-out error by \ph{XX\%} relative to independent memories and improves on the purely additive prior ($\kappa_\times=0$) when surface and load interact. In the disconnected control, the additive prior cannot determine the held-out combination and performs close to independent memories (\cref{fig:results}d), as \cref{prop:composition} predicts.

\paragraph{Calibration, safety, probing and hardware deployment.}
Predictive weights keep the intervals of \cref{thm:confidence} calibrated, unlike filtered weights. At a matched latent violation rate of \ph{2\%}, per-memory margins intervene \ph{XX\%} less often than a single adaptive margin, and value-of-information probing needs \ph{XX\%} fewer probes (\cref{app:additional}). On hardware, \method created \ph{11} memories for \ph{9} conditions, and recall accounted for \ph{XX\%} of the improvement on returns (\cref{app:deployment}).
